% Options for packages loaded elsewhere
\PassOptionsToPackage{unicode}{hyperref}
\PassOptionsToPackage{hyphens}{url}
\documentclass[
]{article}
\usepackage{xcolor}
\usepackage[margin=1in]{geometry}
\usepackage{amsmath,amssymb}
\setcounter{secnumdepth}{-\maxdimen} % remove section numbering
\usepackage{iftex}
\ifPDFTeX
  \usepackage[T1]{fontenc}
  \usepackage[utf8]{inputenc}
  \usepackage{textcomp} % provide euro and other symbols
\else % if luatex or xetex
  \usepackage{unicode-math} % this also loads fontspec
  \defaultfontfeatures{Scale=MatchLowercase}
  \defaultfontfeatures[\rmfamily]{Ligatures=TeX,Scale=1}
\fi
\usepackage{lmodern}
\ifPDFTeX\else
  % xetex/luatex font selection
\fi
% Use upquote if available, for straight quotes in verbatim environments
\IfFileExists{upquote.sty}{\usepackage{upquote}}{}
\IfFileExists{microtype.sty}{% use microtype if available
  \usepackage[]{microtype}
  \UseMicrotypeSet[protrusion]{basicmath} % disable protrusion for tt fonts
}{}
\makeatletter
\@ifundefined{KOMAClassName}{% if non-KOMA class
  \IfFileExists{parskip.sty}{%
    \usepackage{parskip}
  }{% else
    \setlength{\parindent}{0pt}
    \setlength{\parskip}{6pt plus 2pt minus 1pt}}
}{% if KOMA class
  \KOMAoptions{parskip=half}}
\makeatother
\usepackage{color}
\usepackage{fancyvrb}
\newcommand{\VerbBar}{|}
\newcommand{\VERB}{\Verb[commandchars=\\\{\}]}
\DefineVerbatimEnvironment{Highlighting}{Verbatim}{commandchars=\\\{\}}
% Add ',fontsize=\small' for more characters per line
\usepackage{framed}
\definecolor{shadecolor}{RGB}{248,248,248}
\newenvironment{Shaded}{\begin{snugshade}}{\end{snugshade}}
\newcommand{\AlertTok}[1]{\textcolor[rgb]{0.94,0.16,0.16}{#1}}
\newcommand{\AnnotationTok}[1]{\textcolor[rgb]{0.56,0.35,0.01}{\textbf{\textit{#1}}}}
\newcommand{\AttributeTok}[1]{\textcolor[rgb]{0.13,0.29,0.53}{#1}}
\newcommand{\BaseNTok}[1]{\textcolor[rgb]{0.00,0.00,0.81}{#1}}
\newcommand{\BuiltInTok}[1]{#1}
\newcommand{\CharTok}[1]{\textcolor[rgb]{0.31,0.60,0.02}{#1}}
\newcommand{\CommentTok}[1]{\textcolor[rgb]{0.56,0.35,0.01}{\textit{#1}}}
\newcommand{\CommentVarTok}[1]{\textcolor[rgb]{0.56,0.35,0.01}{\textbf{\textit{#1}}}}
\newcommand{\ConstantTok}[1]{\textcolor[rgb]{0.56,0.35,0.01}{#1}}
\newcommand{\ControlFlowTok}[1]{\textcolor[rgb]{0.13,0.29,0.53}{\textbf{#1}}}
\newcommand{\DataTypeTok}[1]{\textcolor[rgb]{0.13,0.29,0.53}{#1}}
\newcommand{\DecValTok}[1]{\textcolor[rgb]{0.00,0.00,0.81}{#1}}
\newcommand{\DocumentationTok}[1]{\textcolor[rgb]{0.56,0.35,0.01}{\textbf{\textit{#1}}}}
\newcommand{\ErrorTok}[1]{\textcolor[rgb]{0.64,0.00,0.00}{\textbf{#1}}}
\newcommand{\ExtensionTok}[1]{#1}
\newcommand{\FloatTok}[1]{\textcolor[rgb]{0.00,0.00,0.81}{#1}}
\newcommand{\FunctionTok}[1]{\textcolor[rgb]{0.13,0.29,0.53}{\textbf{#1}}}
\newcommand{\ImportTok}[1]{#1}
\newcommand{\InformationTok}[1]{\textcolor[rgb]{0.56,0.35,0.01}{\textbf{\textit{#1}}}}
\newcommand{\KeywordTok}[1]{\textcolor[rgb]{0.13,0.29,0.53}{\textbf{#1}}}
\newcommand{\NormalTok}[1]{#1}
\newcommand{\OperatorTok}[1]{\textcolor[rgb]{0.81,0.36,0.00}{\textbf{#1}}}
\newcommand{\OtherTok}[1]{\textcolor[rgb]{0.56,0.35,0.01}{#1}}
\newcommand{\PreprocessorTok}[1]{\textcolor[rgb]{0.56,0.35,0.01}{\textit{#1}}}
\newcommand{\RegionMarkerTok}[1]{#1}
\newcommand{\SpecialCharTok}[1]{\textcolor[rgb]{0.81,0.36,0.00}{\textbf{#1}}}
\newcommand{\SpecialStringTok}[1]{\textcolor[rgb]{0.31,0.60,0.02}{#1}}
\newcommand{\StringTok}[1]{\textcolor[rgb]{0.31,0.60,0.02}{#1}}
\newcommand{\VariableTok}[1]{\textcolor[rgb]{0.00,0.00,0.00}{#1}}
\newcommand{\VerbatimStringTok}[1]{\textcolor[rgb]{0.31,0.60,0.02}{#1}}
\newcommand{\WarningTok}[1]{\textcolor[rgb]{0.56,0.35,0.01}{\textbf{\textit{#1}}}}
\usepackage{longtable,booktabs,array}
\newcounter{none} % for unnumbered tables
\usepackage{calc} % for calculating minipage widths
% Correct order of tables after \paragraph or \subparagraph
\usepackage{etoolbox}
\makeatletter
\patchcmd\longtable{\par}{\if@noskipsec\mbox{}\fi\par}{}{}
\makeatother
% Allow footnotes in longtable head/foot
\IfFileExists{footnotehyper.sty}{\usepackage{footnotehyper}}{\usepackage{footnote}}
\makesavenoteenv{longtable}
\usepackage{graphicx}
\makeatletter
\newsavebox\pandoc@box
\newcommand*\pandocbounded[1]{% scales image to fit in text height/width
  \sbox\pandoc@box{#1}%
  \Gscale@div\@tempa{\textheight}{\dimexpr\ht\pandoc@box+\dp\pandoc@box\relax}%
  \Gscale@div\@tempb{\linewidth}{\wd\pandoc@box}%
  \ifdim\@tempb\p@<\@tempa\p@\let\@tempa\@tempb\fi% select the smaller of both
  \ifdim\@tempa\p@<\p@\scalebox{\@tempa}{\usebox\pandoc@box}%
  \else\usebox{\pandoc@box}%
  \fi%
}
% Set default figure placement to htbp
\def\fps@figure{htbp}
\makeatother
\setlength{\emergencystretch}{3em} % prevent overfull lines
\providecommand{\tightlist}{%
  \setlength{\itemsep}{0pt}\setlength{\parskip}{0pt}}
\usepackage{bookmark}
\IfFileExists{xurl.sty}{\usepackage{xurl}}{} % add URL line breaks if available
\urlstyle{same}
\hypersetup{
  pdftitle={Question 2 Analysis},
  pdfauthor={Siva Lakshmi},
  hidelinks,
  pdfcreator={LaTeX via pandoc}}

\title{Question 2 Analysis}
\author{Siva Lakshmi}
\date{2026-04-24}

\begin{document}
\maketitle

\begin{Shaded}
\begin{Highlighting}[]
\FunctionTok{library}\NormalTok{(tidyverse)}

\NormalTok{raw\_files }\OtherTok{\textless{}{-}} \FunctionTok{c}\NormalTok{(}
  \StringTok{"2019"} \OtherTok{=} \StringTok{"../../data/raw/PSI\_IPC\_2019.csv"}\NormalTok{,}
  \StringTok{"2020"} \OtherTok{=} \StringTok{"../../data/raw/PSI\_IPC\_2020.csv"}\NormalTok{,}
  \StringTok{"2021"} \OtherTok{=} \StringTok{"../../data/raw/PSI\_IPC\_2021.csv"}\NormalTok{,}
  \StringTok{"2022"} \OtherTok{=} \StringTok{"../../data/raw/PSI\_IPC\_2022.csv"}\NormalTok{,}
  \StringTok{"2023"} \OtherTok{=} \StringTok{"../../data/raw/PSI\_IPC\_2023.csv"}
\NormalTok{)}

\NormalTok{process\_file }\OtherTok{\textless{}{-}} \ControlFlowTok{function}\NormalTok{(year, path) \{}
\NormalTok{  df }\OtherTok{\textless{}{-}} \FunctionTok{read.csv}\NormalTok{(path, }\AttributeTok{check.names =} \ConstantTok{FALSE}\NormalTok{, }\AttributeTok{stringsAsFactors =} \ConstantTok{FALSE}\NormalTok{)}
  
\NormalTok{  clean\_df }\OtherTok{\textless{}{-}} \FunctionTok{data.frame}\NormalTok{(}
    \AttributeTok{Year =} \FunctionTok{as.numeric}\NormalTok{(year),}
    \AttributeTok{State =} \FunctionTok{toupper}\NormalTok{(}\FunctionTok{trimws}\NormalTok{(}\FunctionTok{as.character}\NormalTok{(df[[}\DecValTok{2}\NormalTok{]]))), }\CommentTok{\# Col 2: State}
    \AttributeTok{Murder =} \FunctionTok{suppressWarnings}\NormalTok{(}\FunctionTok{as.numeric}\NormalTok{(df[[}\DecValTok{3}\NormalTok{]])), }\CommentTok{\# Col 3: Murder}
    \StringTok{\textasciigrave{}}\AttributeTok{Attempt to commit Murder}\StringTok{\textasciigrave{}} \OtherTok{=} \FunctionTok{suppressWarnings}\NormalTok{(}\FunctionTok{as.numeric}\NormalTok{(df[[}\DecValTok{6}\NormalTok{]])), }\CommentTok{\# Col 6: Attempt}
    \StringTok{\textasciigrave{}}\AttributeTok{Kidnapping \& Abduction}\StringTok{\textasciigrave{}} \OtherTok{=} \FunctionTok{suppressWarnings}\NormalTok{(}\FunctionTok{as.numeric}\NormalTok{(df[[}\DecValTok{7}\NormalTok{]])), }\CommentTok{\# Col 7: Kidnapping}
    \AttributeTok{Theft =} \FunctionTok{suppressWarnings}\NormalTok{(}\FunctionTok{as.numeric}\NormalTok{(df[[}\DecValTok{12}\NormalTok{]])), }\CommentTok{\# Col 12: Theft}
    \AttributeTok{Robbery =} \FunctionTok{suppressWarnings}\NormalTok{(}\FunctionTok{as.numeric}\NormalTok{(df[[}\DecValTok{14}\NormalTok{]])), }\CommentTok{\# Col 14: Robbery}
    \AttributeTok{Burglary =} \FunctionTok{suppressWarnings}\NormalTok{(}\FunctionTok{as.numeric}\NormalTok{(df[[}\DecValTok{20}\NormalTok{]])), }\CommentTok{\# Col 20: Burglary}
    \AttributeTok{check.names =} \ConstantTok{FALSE}
\NormalTok{  )}
  
\NormalTok{  clean\_df }\OtherTok{\textless{}{-}}\NormalTok{ clean\_df }\SpecialCharTok{\%\textgreater{}\%}
    \FunctionTok{filter}\NormalTok{(}\SpecialCharTok{!}\FunctionTok{grepl}\NormalTok{(}\StringTok{"TOTAL|INDIA"}\NormalTok{, State, }\AttributeTok{ignore.case =} \ConstantTok{TRUE}\NormalTok{)) }\SpecialCharTok{\%\textgreater{}\%}
    \FunctionTok{replace}\NormalTok{(}\FunctionTok{is.na}\NormalTok{(.), }\DecValTok{0}\NormalTok{)}
  
  \FunctionTok{return}\NormalTok{(clean\_df)}
\NormalTok{\}}

\NormalTok{crime\_data }\OtherTok{\textless{}{-}} \FunctionTok{map2\_dfr}\NormalTok{(}\FunctionTok{names}\NormalTok{(raw\_files), raw\_files, process\_file)}
\end{Highlighting}
\end{Shaded}

\section{Comprehensive Report: State Crime Profiles and Conviction
Patterns
(2019--2023)}\label{comprehensive-report-state-crime-profiles-and-conviction-patterns-20192023}

\subsection{1. Parent Question}\label{parent-question}

\textbf{Do states exhibit distinct ``crime profiles'' based on the
composition of convictions?} \emph{(Note: This overarching inquiry seeks
to classify states based on their criminal conviction outputs, analyzing
shifts between violence and property crimes, identifying lethal
outliers, and tracking the volatility of specific offenses to inform
state-level prison administration and police deployment.)}

\begin{center}\rule{0.5\linewidth}{0.5pt}\end{center}

\subsection{Sub-Question 2.1: Classifying Prison Populations (Violence
vs.~Property)}\label{sub-question-2.1-classifying-prison-populations-violence-vs.-property}

\textbf{Using a simple ratio metric, can we classify states on a
spectrum from ``Violence-Dominant'' to ``Property-Dominant'' prisons?
How has Uttar Pradesh's position on this spectrum shifted from 2019 to
2023 compared to Tamil Nadu?}

\subsubsection{Statistical Methods Used}\label{statistical-methods-used}

\begin{itemize}
\tightlist
\item
  Composite Index Calculation
\item
  Comparative Time Series Plot Visualization
\end{itemize}

\subsubsection{Step-by-Step Application}\label{step-by-step-application}

\begin{enumerate}
\def\labelenumi{\arabic{enumi}.}
\tightlist
\item
  \textbf{Define Ratio:} Created a \texttt{Violence\_Index} calculated
  as the sum of violent crimes (Murder, Attempt to Murder, Kidnapping)
  divided by property crimes (Theft, Burglary, Robbery).
\item
  \textbf{Calculate Yearly:} Computed this index for every state across
  the 5-year period (2019--2023).
\item
  \textbf{National Benchmarking:} Calculated the national average for
  the index across the same period.
\item
  \textbf{Visualization:} Plotted a time series line chart isolating
  Uttar Pradesh (UP), Tamil Nadu (TN), and the National Average to
  observe relative shifts.
\end{enumerate}

\subsubsection{R Code}\label{r-code}

\begin{Shaded}
\begin{Highlighting}[]
\FunctionTok{library}\NormalTok{(tidyverse)}

\NormalTok{index\_data }\OtherTok{\textless{}{-}}\NormalTok{ crime\_data }\SpecialCharTok{\%\textgreater{}\%}
  \FunctionTok{mutate}\NormalTok{(}
    \AttributeTok{Violence\_Index =}\NormalTok{ (Murder }\SpecialCharTok{+} \StringTok{\textasciigrave{}}\AttributeTok{Attempt to commit Murder}\StringTok{\textasciigrave{}} \SpecialCharTok{+} \StringTok{\textasciigrave{}}\AttributeTok{Kidnapping \& Abduction}\StringTok{\textasciigrave{}}\NormalTok{) }\SpecialCharTok{/} 
\NormalTok{                     (Theft }\SpecialCharTok{+}\NormalTok{ Burglary }\SpecialCharTok{+}\NormalTok{ Robbery }\SpecialCharTok{+} \FloatTok{0.0001}\NormalTok{)}
\NormalTok{  )}

\NormalTok{national\_avg }\OtherTok{\textless{}{-}}\NormalTok{ index\_data }\SpecialCharTok{\%\textgreater{}\%}
  \FunctionTok{group\_by}\NormalTok{(Year) }\SpecialCharTok{\%\textgreater{}\%}
  \FunctionTok{summarize}\NormalTok{(}\AttributeTok{Violence\_Index =} \FunctionTok{mean}\NormalTok{(Violence\_Index, }\AttributeTok{na.rm =} \ConstantTok{TRUE}\NormalTok{)) }\SpecialCharTok{\%\textgreater{}\%}
  \FunctionTok{mutate}\NormalTok{(}\AttributeTok{State =} \StringTok{"NATIONAL AVERAGE"}\NormalTok{)}

\NormalTok{plot\_data\_q1 }\OtherTok{\textless{}{-}}\NormalTok{ index\_data }\SpecialCharTok{\%\textgreater{}\%}
  \FunctionTok{filter}\NormalTok{(}\FunctionTok{toupper}\NormalTok{(}\FunctionTok{trimws}\NormalTok{(State)) }\SpecialCharTok{\%in\%} \FunctionTok{c}\NormalTok{(}\StringTok{"UTTAR PRADESH"}\NormalTok{, }\StringTok{"TAMIL NADU"}\NormalTok{)) }\SpecialCharTok{\%\textgreater{}\%}
  \FunctionTok{select}\NormalTok{(Year, State, Violence\_Index) }\SpecialCharTok{\%\textgreater{}\%}
  \FunctionTok{bind\_rows}\NormalTok{(national\_avg)}

\FunctionTok{ggplot}\NormalTok{(plot\_data\_q1, }\FunctionTok{aes}\NormalTok{(}\AttributeTok{x =}\NormalTok{ Year, }\AttributeTok{y =}\NormalTok{ Violence\_Index, }\AttributeTok{color =}\NormalTok{ State, }\AttributeTok{group =}\NormalTok{ State)) }\SpecialCharTok{+}
  \FunctionTok{geom\_line}\NormalTok{(}\AttributeTok{linewidth =} \FloatTok{1.2}\NormalTok{) }\SpecialCharTok{+}
  \FunctionTok{geom\_point}\NormalTok{(}\AttributeTok{size =} \DecValTok{3}\NormalTok{) }\SpecialCharTok{+}
  \FunctionTok{scale\_color\_brewer}\NormalTok{(}\AttributeTok{palette =} \StringTok{"Set1"}\NormalTok{) }\SpecialCharTok{+}
  \FunctionTok{labs}\NormalTok{(}\AttributeTok{title =} \StringTok{"Violence Index: UP vs Tamil Nadu vs National Avg (2019{-}2023)"}\NormalTok{,}
       \AttributeTok{y =} \StringTok{"Violence Index (Higher = More Violent)"}\NormalTok{, }\AttributeTok{x =} \StringTok{"Year"}\NormalTok{, }\AttributeTok{color =} \StringTok{"State/Avg"}\NormalTok{) }\SpecialCharTok{+}
  \FunctionTok{theme\_minimal}\NormalTok{()}
\end{Highlighting}
\end{Shaded}

\pandocbounded{\includegraphics[keepaspectratio]{q2_analysis_files/figure-latex/unnamed-chunk-1-1.pdf}}

\subsubsection{Interpretation}\label{interpretation}

\begin{itemize}
\tightlist
\item
  \textbf{Distinct Regional Profiles:} * *
\item
  \textbf{Shifts in Uttar Pradesh:} * *
\item
  \textbf{Conclusion:} The data confirms the existence of distinct
  regional crime profiles. States on the higher end of this spectrum
  require prison infrastructure focused on maximum security and mental
  health intervention, whereas states on the lower end must prioritize
  vocational training for property offenders.
\end{itemize}

\begin{center}\rule{0.5\linewidth}{0.5pt}\end{center}

\subsection{Sub-Question 2.2: Rank Mobility in Murder
Convictions}\label{sub-question-2.2-rank-mobility-in-murder-convictions}

\textbf{How many states changed their rank position in the Murder
Conviction Count by more than 3 places between 2019 and 2023? Is the
``Murder Capital'' list rigid or fluid?}

\subsubsection{Statistical Methods
Used}\label{statistical-methods-used-1}

\begin{itemize}
\tightlist
\item
  Rank Mobility Index
\item
  Tabular Comparison
\end{itemize}

\subsubsection{Step-by-Step
Application}\label{step-by-step-application-1}

\begin{enumerate}
\def\labelenumi{\arabic{enumi}.}
\tightlist
\item
  \textbf{Ranking 2019:} Ranked all states based on their total Murder
  convictions in 2019.
\item
  \textbf{Ranking 2023:} Ranked all states based on their total Murder
  convictions in 2023.
\item
  \textbf{Calculate Shift:} Joined the datasets and calculated the
  absolute difference between the two ranks.
\item
  \textbf{Threshold Filtering:} Isolated states that experienced a rank
  shift of more than 3 positions.
\end{enumerate}

\subsubsection{R Code}\label{r-code-1}

\begin{Shaded}
\begin{Highlighting}[]
\NormalTok{rank\_2019 }\OtherTok{\textless{}{-}}\NormalTok{ crime\_data }\SpecialCharTok{\%\textgreater{}\%}
  \FunctionTok{filter}\NormalTok{(Year }\SpecialCharTok{==} \DecValTok{2019}\NormalTok{) }\SpecialCharTok{\%\textgreater{}\%}
  \FunctionTok{mutate}\NormalTok{(}\AttributeTok{Rank\_2019 =} \FunctionTok{min\_rank}\NormalTok{(}\FunctionTok{desc}\NormalTok{(Murder))) }\SpecialCharTok{\%\textgreater{}\%}
  \FunctionTok{select}\NormalTok{(State, Rank\_2019)}

\NormalTok{rank\_2023 }\OtherTok{\textless{}{-}}\NormalTok{ crime\_data }\SpecialCharTok{\%\textgreater{}\%}
  \FunctionTok{filter}\NormalTok{(Year }\SpecialCharTok{==} \DecValTok{2023}\NormalTok{) }\SpecialCharTok{\%\textgreater{}\%}
  \FunctionTok{mutate}\NormalTok{(}\AttributeTok{Rank\_2023 =} \FunctionTok{min\_rank}\NormalTok{(}\FunctionTok{desc}\NormalTok{(Murder))) }\SpecialCharTok{\%\textgreater{}\%}
  \FunctionTok{select}\NormalTok{(State, Rank\_2023)}

\NormalTok{rank\_mobility }\OtherTok{\textless{}{-}} \FunctionTok{inner\_join}\NormalTok{(rank\_2019, rank\_2023, }\AttributeTok{by =} \StringTok{"State"}\NormalTok{) }\SpecialCharTok{\%\textgreater{}\%}
  \FunctionTok{mutate}\NormalTok{(}\AttributeTok{Rank\_Change =}\NormalTok{ Rank\_2019 }\SpecialCharTok{{-}}\NormalTok{ Rank\_2023) }\SpecialCharTok{\%\textgreater{}\%} 
  \FunctionTok{filter}\NormalTok{(}\FunctionTok{abs}\NormalTok{(Rank\_Change) }\SpecialCharTok{\textgreater{}} \DecValTok{3}\NormalTok{) }\SpecialCharTok{\%\textgreater{}\%}
  \FunctionTok{arrange}\NormalTok{(}\FunctionTok{desc}\NormalTok{(Rank\_Change))}

\NormalTok{knitr}\SpecialCharTok{::}\FunctionTok{kable}\NormalTok{(rank\_mobility, }\AttributeTok{caption =} \StringTok{"States with Rank Shifts \textgreater{} 3 Places (Murder Convictions)"}\NormalTok{)}
\end{Highlighting}
\end{Shaded}

\begin{longtable}[]{@{}lrrr@{}}
\caption{States with Rank Shifts \textgreater{} 3 Places (Murder
Convictions)}\tabularnewline
\toprule\noalign{}
State & Rank\_2019 & Rank\_2023 & Rank\_Change \\
\midrule\noalign{}
\endfirsthead
\toprule\noalign{}
State & Rank\_2019 & Rank\_2023 & Rank\_Change \\
\midrule\noalign{}
\endhead
\bottomrule\noalign{}
\endlastfoot
GUJARAT & 13 & 6 & 7 \\
ARUNACHAL PRADESH & 33 & 27 & 6 \\
GOA & 25 & 29 & -4 \\
HARYANA & 6 & 10 & -4 \\
KARNATAKA & 10 & 14 & -4 \\
\end{longtable}

\subsubsection{Interpretation}\label{interpretation-1}

\begin{itemize}
\tightlist
\item
  \textbf{List Fluidity:} * *
\item
  \textbf{Structural Inertia:} * *
\item
  \textbf{Conclusion:} Significant rank climbs serve as a red alert for
  policy makers, highlighting areas where new, non-traditional violence
  has begun entering the court system.
\end{itemize}

\begin{center}\rule{0.5\linewidth}{0.5pt}\end{center}

\subsection{Sub-Question 2.3: Assessing Lethality (Attempted vs.~Actual
Murder)}\label{sub-question-2.3-assessing-lethality-attempted-vs.-actual-murder}

\textbf{Which states have an unusually low ratio of Attempt to Murder to
Murder compared to their regional neighbors?}

\subsubsection{Statistical Methods
Used}\label{statistical-methods-used-2}

\begin{itemize}
\tightlist
\item
  Ratio Analysis
\item
  Outlier Box Plot Visualization
\end{itemize}

\subsubsection{Step-by-Step
Application}\label{step-by-step-application-2}

\begin{enumerate}
\def\labelenumi{\arabic{enumi}.}
\tightlist
\item
  \textbf{Lethality Calculation:} Computed the average ratio of Attempt
  to Murder versus actual Murder convictions for each state.
\item
  \textbf{Outlier Detection:} Utilized the Interquartile Range (IQR)
  method to identify statistical lower outliers.
\item
  \textbf{Visualization:} Plotted the distribution on a box plot,
  explicitly labeling states that fall below the normal distribution
  threshold.
\end{enumerate}

\subsubsection{R Code}\label{r-code-2}

\begin{Shaded}
\begin{Highlighting}[]
\NormalTok{lethality\_data }\OtherTok{\textless{}{-}}\NormalTok{ crime\_data }\SpecialCharTok{\%\textgreater{}\%}
  \FunctionTok{group\_by}\NormalTok{(State) }\SpecialCharTok{\%\textgreater{}\%}
  \FunctionTok{summarise}\NormalTok{(}
    \AttributeTok{Avg\_Murder =} \FunctionTok{mean}\NormalTok{(Murder, }\AttributeTok{na.rm =} \ConstantTok{TRUE}\NormalTok{),}
    \AttributeTok{Avg\_Attempt =} \FunctionTok{mean}\NormalTok{(}\StringTok{\textasciigrave{}}\AttributeTok{Attempt to commit Murder}\StringTok{\textasciigrave{}}\NormalTok{, }\AttributeTok{na.rm =} \ConstantTok{TRUE}\NormalTok{)}
\NormalTok{  ) }\SpecialCharTok{\%\textgreater{}\%}
  \FunctionTok{mutate}\NormalTok{(}\AttributeTok{Lethality\_Ratio =}\NormalTok{ Avg\_Attempt }\SpecialCharTok{/}\NormalTok{ (Avg\_Murder }\SpecialCharTok{+} \FloatTok{0.0001}\NormalTok{))}

\NormalTok{q1 }\OtherTok{\textless{}{-}} \FunctionTok{quantile}\NormalTok{(lethality\_data}\SpecialCharTok{$}\NormalTok{Lethality\_Ratio, }\FloatTok{0.25}\NormalTok{, }\AttributeTok{na.rm =} \ConstantTok{TRUE}\NormalTok{)}
\NormalTok{iqr\_val }\OtherTok{\textless{}{-}} \FunctionTok{IQR}\NormalTok{(lethality\_data}\SpecialCharTok{$}\NormalTok{Lethality\_Ratio, }\AttributeTok{na.rm =} \ConstantTok{TRUE}\NormalTok{)}
\NormalTok{lower\_bound }\OtherTok{\textless{}{-}}\NormalTok{ q1 }\SpecialCharTok{{-}} \FloatTok{1.5} \SpecialCharTok{*}\NormalTok{ iqr\_val}

\FunctionTok{ggplot}\NormalTok{(lethality\_data, }\FunctionTok{aes}\NormalTok{(}\AttributeTok{y =}\NormalTok{ Lethality\_Ratio)) }\SpecialCharTok{+}
  \FunctionTok{geom\_boxplot}\NormalTok{(}\AttributeTok{fill =} \StringTok{"lightblue"}\NormalTok{, }\AttributeTok{alpha =} \FloatTok{0.7}\NormalTok{) }\SpecialCharTok{+}
  \FunctionTok{geom\_text}\NormalTok{(}\FunctionTok{aes}\NormalTok{(}\AttributeTok{x =} \DecValTok{0}\NormalTok{, }\AttributeTok{label =} \FunctionTok{ifelse}\NormalTok{(Lethality\_Ratio }\SpecialCharTok{\textless{}}\NormalTok{ lower\_bound, State, }\StringTok{""}\NormalTok{)), }
            \AttributeTok{hjust =} \SpecialCharTok{{-}}\FloatTok{0.5}\NormalTok{, }\AttributeTok{color =} \StringTok{"red"}\NormalTok{, }\AttributeTok{fontface =} \StringTok{"bold"}\NormalTok{) }\SpecialCharTok{+}
  \FunctionTok{labs}\NormalTok{(}\AttributeTok{title =} \StringTok{"Lethality: Attempt to Murder vs Murder Ratio"}\NormalTok{,}
       \AttributeTok{y =} \StringTok{"Ratio (Attempts per 1 Murder)"}\NormalTok{) }\SpecialCharTok{+}
  \FunctionTok{theme\_minimal}\NormalTok{() }\SpecialCharTok{+}
  \FunctionTok{theme}\NormalTok{(}\AttributeTok{axis.text.x =} \FunctionTok{element\_blank}\NormalTok{(), }\AttributeTok{axis.ticks.x =} \FunctionTok{element\_blank}\NormalTok{())}
\end{Highlighting}
\end{Shaded}

\pandocbounded{\includegraphics[keepaspectratio]{q2_analysis_files/figure-latex/unnamed-chunk-3-1.pdf}}

\subsubsection{Interpretation}\label{interpretation-2}

\begin{itemize}
\tightlist
\item
  \textbf{High Lethality Outliers:} * *
\item
  \textbf{Forensic Evidence of Infrastructure:} This low ratio serves as
  a proxy indicating either widespread availability of lethal weaponry
  or systemic failures in emergency trauma care.
\item
  \textbf{Conclusion:} States identified as lower outliers require
  immediate public health and policing interventions, specifically
  focusing on weapon confiscation and trauma center accessibility.
\end{itemize}

\begin{center}\rule{0.5\linewidth}{0.5pt}\end{center}

\subsection{Sub-Question 2.4: Volatility in Kidnapping and Abduction
Convictions}\label{sub-question-2.4-volatility-in-kidnapping-and-abduction-convictions}

\textbf{Group states into Low, Medium, and High tertiles based on their
5-year average Kidnapping \& Abduction convictions. Which states in the
Medium tertile show High Volatility (CV \textgreater{} 0.5)?}

\subsubsection{Statistical Methods
Used}\label{statistical-methods-used-3}

\begin{itemize}
\tightlist
\item
  Tertile Grouping (Rank-Based Classification)
\item
  Coefficient of Variation (CV = SD / Mean)
\item
  Threshold Filtering (CV \textgreater{} 0.5)
\end{itemize}

\subsubsection{Step-by-Step
Application}\label{step-by-step-application-3}

\begin{enumerate}
\def\labelenumi{\arabic{enumi}.}
\tightlist
\item
  \textbf{Tertile Assignment:} Calculated the 5-year mean for kidnapping
  convictions and split the states into three equal tertiles.
\item
  \textbf{Volatility Calculation:} Isolated the ``Medium'' tertile and
  computed the Standard Deviation and Coefficient of Variation (CV) for
  each state.
\item
  \textbf{Threshold Isolation:} Filtered out states with a CV greater
  than 0.5 (indicating standard deviation is half the mean).
\item
  \textbf{Visualization:} Plotted the 5-year trend for these highly
  volatile states to visualize the erratic nature of the convictions.
\end{enumerate}

\subsubsection{R Code}\label{r-code-3}

\begin{Shaded}
\begin{Highlighting}[]
\NormalTok{kidnap\_stats }\OtherTok{\textless{}{-}}\NormalTok{ crime\_data }\SpecialCharTok{\%\textgreater{}\%}
  \FunctionTok{group\_by}\NormalTok{(State) }\SpecialCharTok{\%\textgreater{}\%}
  \FunctionTok{summarise}\NormalTok{(}
    \AttributeTok{Mean\_Kidnap =} \FunctionTok{mean}\NormalTok{(}\StringTok{\textasciigrave{}}\AttributeTok{Kidnapping \& Abduction}\StringTok{\textasciigrave{}}\NormalTok{, }\AttributeTok{na.rm =} \ConstantTok{TRUE}\NormalTok{),}
    \AttributeTok{SD\_Kidnap =} \FunctionTok{sd}\NormalTok{(}\StringTok{\textasciigrave{}}\AttributeTok{Kidnapping \& Abduction}\StringTok{\textasciigrave{}}\NormalTok{, }\AttributeTok{na.rm =} \ConstantTok{TRUE}\NormalTok{)}
\NormalTok{  ) }\SpecialCharTok{\%\textgreater{}\%}
  \FunctionTok{mutate}\NormalTok{(}
    \AttributeTok{CV =}\NormalTok{ SD\_Kidnap }\SpecialCharTok{/}\NormalTok{ (Mean\_Kidnap }\SpecialCharTok{+} \FloatTok{0.0001}\NormalTok{),}
    \AttributeTok{Tertile =} \FunctionTok{ntile}\NormalTok{(Mean\_Kidnap, }\DecValTok{3}\NormalTok{) }
\NormalTok{  )}

\NormalTok{volatile\_medium\_states }\OtherTok{\textless{}{-}}\NormalTok{ kidnap\_stats }\SpecialCharTok{\%\textgreater{}\%}
  \FunctionTok{filter}\NormalTok{(Tertile }\SpecialCharTok{==} \DecValTok{2}\NormalTok{, CV }\SpecialCharTok{\textgreater{}} \FloatTok{0.5}\NormalTok{) }\SpecialCharTok{\%\textgreater{}\%}
  \FunctionTok{pull}\NormalTok{(State)}

\NormalTok{volatile\_data }\OtherTok{\textless{}{-}}\NormalTok{ crime\_data }\SpecialCharTok{\%\textgreater{}\%}
  \FunctionTok{filter}\NormalTok{(State }\SpecialCharTok{\%in\%}\NormalTok{ volatile\_medium\_states)}

\ControlFlowTok{if}\NormalTok{(}\FunctionTok{length}\NormalTok{(volatile\_medium\_states) }\SpecialCharTok{\textgreater{}} \DecValTok{0}\NormalTok{) \{}
  \FunctionTok{ggplot}\NormalTok{(volatile\_data, }\FunctionTok{aes}\NormalTok{(}\AttributeTok{x =}\NormalTok{ Year, }\AttributeTok{y =} \StringTok{\textasciigrave{}}\AttributeTok{Kidnapping \& Abduction}\StringTok{\textasciigrave{}}\NormalTok{, }\AttributeTok{color =}\NormalTok{ State)) }\SpecialCharTok{+}
    \FunctionTok{geom\_line}\NormalTok{(}\AttributeTok{linewidth =} \DecValTok{1}\NormalTok{) }\SpecialCharTok{+}
    \FunctionTok{geom\_point}\NormalTok{(}\AttributeTok{size =} \DecValTok{2}\NormalTok{) }\SpecialCharTok{+}
    \FunctionTok{scale\_color\_brewer}\NormalTok{(}\AttributeTok{palette =} \StringTok{"Dark2"}\NormalTok{) }\SpecialCharTok{+}
    \FunctionTok{labs}\NormalTok{(}\AttributeTok{title =} \StringTok{"High Volatility Kidnapping Trends (Medium Tertile States)"}\NormalTok{,}
         \AttributeTok{y =} \StringTok{"Kidnapping Convictions"}\NormalTok{, }\AttributeTok{x =} \StringTok{"Year"}\NormalTok{) }\SpecialCharTok{+}
    \FunctionTok{theme\_minimal}\NormalTok{()}
\NormalTok{\} }\ControlFlowTok{else}\NormalTok{ \{}
  \FunctionTok{cat}\NormalTok{(}\StringTok{"No states in the Medium Tertile exhibited High Volatility (CV \textgreater{} 0.5).}\SpecialCharTok{\textbackslash{}n}\StringTok{"}\NormalTok{)}
\NormalTok{\}}
\end{Highlighting}
\end{Shaded}

\pandocbounded{\includegraphics[keepaspectratio]{q2_analysis_files/figure-latex/unnamed-chunk-4-1.pdf}}

\subsubsection{Interpretation}\label{interpretation-3}

\begin{itemize}
\tightlist
\item
  \textbf{The ``Invisible Middle'':} * *
\item
  \textbf{Resource Misallocation:} Because their average numbers look
  stable, these states likely lack specialized Anti-Human Trafficking
  Units (AHTUs) and are unprepared for peak surge years.
\item
  \textbf{Conclusion:} States identified with high CVs require flexible,
  rapid-response task forces rather than standard permanent sizing, as
  their kidnapping rates are episodic rather than endemic.
\end{itemize}

\begin{center}\rule{0.5\linewidth}{0.5pt}\end{center}

\subsection{Sub-Question 2.5: Correlation Between Robbery and
Burglary}\label{sub-question-2.5-correlation-between-robbery-and-burglary}

\textbf{For each state, calculate the 5-year correlation (Pearson's r)
between Robbery and Burglary. Which states exhibit a strong negative
correlation (r \textless{} -0.7)?}

\subsubsection{Statistical Methods
Used}\label{statistical-methods-used-4}

\begin{itemize}
\tightlist
\item
  Pearson Correlation Coefficient (r)
\item
  Scatter Plot with Linear Trendline
\end{itemize}

\subsubsection{Step-by-Step
Application}\label{step-by-step-application-4}

\begin{enumerate}
\def\labelenumi{\arabic{enumi}.}
\tightlist
\item
  \textbf{Correlation Calculation:} Grouped data by state and computed
  Pearson's r between Robbery and Burglary across the 5 paired annual
  observations.
\item
  \textbf{Threshold Filtering:} Identified states demonstrating a strong
  negative correlation (r \textless{} -0.7).
\item
  \textbf{Visualization:} Generated scatter plots for the flagged
  states, mapping Robbery on the X-axis and Burglary on the Y-axis,
  overlaid with a linear regression line.
\end{enumerate}

\subsubsection{R Code}\label{r-code-4}

\begin{Shaded}
\begin{Highlighting}[]
\NormalTok{cor\_data }\OtherTok{\textless{}{-}}\NormalTok{ crime\_data }\SpecialCharTok{\%\textgreater{}\%}
  \FunctionTok{group\_by}\NormalTok{(State) }\SpecialCharTok{\%\textgreater{}\%}
  \FunctionTok{summarise}\NormalTok{(}
    \AttributeTok{Correlation =} \FunctionTok{cor}\NormalTok{(Robbery, Burglary, }\AttributeTok{use =} \StringTok{"pairwise.complete.obs"}\NormalTok{)}
\NormalTok{  ) }\SpecialCharTok{\%\textgreater{}\%}
  \FunctionTok{filter}\NormalTok{(Correlation }\SpecialCharTok{\textless{}} \SpecialCharTok{{-}}\FloatTok{0.7}\NormalTok{)}
\end{Highlighting}
\end{Shaded}

\begin{verbatim}
## Warning: There were 10 warnings in `summarise()`.
## The first warning was:
## i In argument: `Correlation = cor(Robbery, Burglary, use =
##   "pairwise.complete.obs")`.
## i In group 1: `State = "A & N ISLANDS"`.
## Caused by warning in `cor()`:
## ! the standard deviation is zero
## i Run `dplyr::last_dplyr_warnings()` to see the 9 remaining warnings.
\end{verbatim}

\begin{Shaded}
\begin{Highlighting}[]
\NormalTok{negative\_cor\_states }\OtherTok{\textless{}{-}}\NormalTok{ cor\_data }\SpecialCharTok{\%\textgreater{}\%} \FunctionTok{pull}\NormalTok{(State)}

\NormalTok{plot\_cor\_data }\OtherTok{\textless{}{-}}\NormalTok{ crime\_data }\SpecialCharTok{\%\textgreater{}\%}
  \FunctionTok{filter}\NormalTok{(State }\SpecialCharTok{\%in\%}\NormalTok{ negative\_cor\_states)}

\ControlFlowTok{if}\NormalTok{(}\FunctionTok{length}\NormalTok{(negative\_cor\_states) }\SpecialCharTok{\textgreater{}} \DecValTok{0}\NormalTok{) \{}
  \FunctionTok{ggplot}\NormalTok{(plot\_cor\_data, }\FunctionTok{aes}\NormalTok{(}\AttributeTok{x =}\NormalTok{ Robbery, }\AttributeTok{y =}\NormalTok{ Burglary)) }\SpecialCharTok{+}
    \FunctionTok{geom\_point}\NormalTok{(}\AttributeTok{size =} \DecValTok{3}\NormalTok{, }\AttributeTok{color =} \StringTok{"darkred"}\NormalTok{) }\SpecialCharTok{+}
    \FunctionTok{geom\_smooth}\NormalTok{(}\AttributeTok{method =} \StringTok{"lm"}\NormalTok{, }\AttributeTok{se =} \ConstantTok{FALSE}\NormalTok{, }\AttributeTok{color =} \StringTok{"steelblue"}\NormalTok{, }\AttributeTok{linewidth =} \DecValTok{1}\NormalTok{) }\SpecialCharTok{+}
    \FunctionTok{facet\_wrap}\NormalTok{(}\SpecialCharTok{\textasciitilde{}}\NormalTok{State, }\AttributeTok{scales =} \StringTok{"free"}\NormalTok{) }\SpecialCharTok{+}
    \FunctionTok{labs}\NormalTok{(}\AttributeTok{title =} \StringTok{"Robbery vs Burglary Substitution (r \textless{} {-}0.7)"}\NormalTok{,}
         \AttributeTok{x =} \StringTok{"Robbery Convictions"}\NormalTok{, }\AttributeTok{y =} \StringTok{"Burglary Convictions"}\NormalTok{) }\SpecialCharTok{+}
    \FunctionTok{theme\_minimal}\NormalTok{()}
\NormalTok{\} }\ControlFlowTok{else}\NormalTok{ \{}
  \FunctionTok{cat}\NormalTok{(}\StringTok{"No states exhibited a strong negative correlation (r \textless{} {-}0.7).}\SpecialCharTok{\textbackslash{}n}\StringTok{"}\NormalTok{)}
\NormalTok{\}}
\end{Highlighting}
\end{Shaded}

\begin{verbatim}
## No states exhibited a strong negative correlation (r < -0.7).
\end{verbatim}

\subsubsection{Interpretation}\label{interpretation-4}

\begin{itemize}
\tightlist
\item
  \textbf{Crime Substitution Effect:} * *
\item
  \textbf{Fluid Criminal Markets:} This strong negative correlation
  demonstrates a highly rational adaptation by offenders to changing
  environmental conditions (e.g., increased home security driving
  criminals to public spaces).
\item
  \textbf{Conclusion:} Police departments in these specific states must
  adopt holistic property crime strategies. Efforts that successfully
  reduce burglary must be immediately paired with transit and
  street-level mitigation measures, or the crime simply relocates to
  robbery.
\end{itemize}

\end{document}
